% ============================================================
%  IV. Proposed RTGS Methodology（提出的 RTGS 方法）
%  中文版本，IEEE 双栏格式
%  符号体系与 Section III (system_description.tex) 保持一致：
%    v_i  = 任务节点      u_k  = 无人机      t = 决策时刻
%    \mathcal{G}_T        = 任务 DAG
%    \mathcal{G}_R        = 资源图
%    \mathbf{f}(t)        = 六维调度反馈向量
%  所有常量均对应代码：
%    configs/params.py  · models/gnn_policy.py  · env/uav_env.py
%    core/task_shaping.py · core/mobility.py    · main.py
% ============================================================

\section{提出的RTGS方法}
\label{sec:method}

本节将资源感知任务图整形问题（RTGS）形式化为马尔可夫决策过程（MDP），
并详细阐述求解该MDP的双图GNN–PPO框架。
在每个决策时刻$t$，智能体观测任务DAG $\mathcal{G}_T$
与资源图$\mathcal{G}_R$的联合编码状态，
输出一个图整形动作，
由确定性贪心调度器完成任务–无人机匹配后，
环境向智能体反馈奖励信号与更新后的状态，
形成如图~\ref{fig:architecture}所示的闭环控制流。

% ----------------------------------------------------------------
\subsection{MDP问题建模}
\label{subsec:mdp}

\subsubsection*{状态空间 $\mathcal{S}$}

系统状态$s_t \in \mathcal{S}$为一个复合观测，
由三个子组件拼合构成
（对应 \texttt{env/uav\_env.py} 中 \texttt{\_get\_obs()} 函数
以及 \texttt{UAVTaskEnv.observation\_space}）。

\begin{enumerate}[label=(\roman*), leftmargin=*, itemsep=3pt]

  \item \textbf{任务DAG观测}：节点特征矩阵
        $\mathbf{X}^T \in \mathbb{R}^{N_{\max} \times 2}$
        与COO格式边索引
        $\mathbf{E}^T \in \mathbb{Z}^{2 \times N_{\max}^2}$，
        节点数上限$N_{\max} = 20$，
        实际节点数$N_t \leq N_{\max}$，
        不足部分以零行填充。
        每行特征$\mathbf{x}_i^T = (c_i, s_i)$与
        第~\ref{sec:system}节式~\eqref{eq:task_attr}中的定义保持一致。

  \item \textbf{资源图观测}：
        由移动性模型动态生成的节点特征矩阵
        $\mathbf{X}^R(t) \in [0,1]^{M \times 3}$
        及全连接边索引$\mathbf{E}^R$；
        第$k$行特征向量的定义见式~\eqref{eq:res_features}：
        \begin{equation}
          \mathbf{x}_k^R(t) =
          \Bigl[\,
            \tilde{\phi}_k(t)/\phi^{\max},\;
            \bar{B}_k(t)/B_{\max},\;
            \beta_k(t)
          \,\Bigr]^\top.
          \label{eq:res_feat_method}
        \end{equation}
        三个分量分别表示
        电量修正后的CPU比率（对应 \texttt{effective\_cpu()} 函数）、
        至活跃邻居的平均链路带宽比以及当前电量，
        均归一化至$[0,1]$。

  \item \textbf{调度反馈向量}：
        六维向量$\mathbf{f}(t) \in \mathbb{R}^6$
        的定义见式~\eqref{eq:feedback}，
        由 \texttt{SchedulingFeedback.as\_vector()} 输出，
        编码了上一轮调度的延迟比$\lambda(t)$、
        任务成功率$\sigma(t)$、强制重调度次数$\varrho(t)$、
        机队平均利用率$\bar{u}(t)$、
        最弱活跃链路带宽比$\hat{B}_{\min}(t)$
        以及最低电量比$\hat{\beta}_{\min}(t)$六个信号。
\end{enumerate}

完整状态可紧凑地表示为
\begin{equation}
  s_t = \bigl(
    \mathbf{X}^T,\;\mathbf{E}^T,\;
    \mathbf{X}^R(t),\;\mathbf{E}^R,\;
    \mathbf{f}(t)
  \bigr).
  \label{eq:state}
\end{equation}

\subsubsection*{动作空间 $\mathcal{A}$}

动作空间为有限离散集（对应 \texttt{configs/params.py} 中常量
\texttt{ACTION\_DIM}）：
\begin{equation}
  |\mathcal{A}| = 1 + N_{\max} + \binom{N_{\max}}{2}
               = 1 + 20 + 190 = 211.
  \label{eq:action_dim}
\end{equation}
三类动作的具体定义如下：
\begin{itemize}[leftmargin=*, itemsep=2pt]
  \item $a = 0$：\textbf{空操作（noop）}，图拓扑保持不变。
  \item $a \in \{1,\ldots,20\}$：
        \textbf{Split$(v_{a-1})$}，
        对填充位置第$a$个节点执行第~\ref{subsec:task_shaping}节
        所定义的\textsc{Split}算子。
  \item $a \in \{21,\ldots,210\}$：
        \textbf{Merge$(v_i, v_j)$}，
        对第$(a-21)$个节点对（按字典序枚举全部$\binom{20}{2}$对）
        执行\textsc{Merge}算子。
\end{itemize}
平坦动作索引到操作类型与节点位置三元组
$(op,\, \text{pad}_i,\, \text{pad}_j)$的解码
由 \texttt{decode\_node\_action()} 函数实现
（\texttt{configs/params.py:81--99}）。

\subsubsection*{奖励函数 $\mathcal{R}$}

奖励函数$r_t$由多个分量叠加而成
（对应 \texttt{env/uav\_env.py:311--336}
中 \texttt{\_compute\_reward()} 函数），
旨在同时激励即时调度质量与增量性能改进：
\begin{equation}
  r_t = r_{\text{qual}} + r_{\Delta} + r_{\text{pen}}.
  \label{eq:reward}
\end{equation}

\noindent\textbf{绝对质量项}：对当前步调度质量给予正向激励，
使智能体在维持良好性能时仍能获得持续奖励：
\begin{equation}
  r_{\text{qual}} = 2\,\sigma(t) + 1.5\,\max\!\bigl(0,\, 1 - \lambda(t)\bigr).
  \label{eq:reward_qual}
\end{equation}

\noindent\textbf{增量改进项}：激励相对于前一决策步的性能提升，
惩罚性能退化：
\begin{equation}
  r_{\Delta} = \bigl(\lambda(t-1) - \lambda(t)\bigr)
             + 0.5\,\bigl(\sigma(t) - \sigma(t-1)\bigr),
  \label{eq:reward_delta}
\end{equation}
延迟比降低（延迟改善）和成功率提升均产生正向增量奖励。

\noindent\textbf{惩罚项}：包含三个负向分量：
\begin{equation}
  r_{\text{pen}} =
    -\frac{\varrho(t)}{N_t}
    - 2\,\max\!\bigl(0,\, \lambda(t)-1\bigr)
    - 0.05\,\max\!\bigl(0,\, N_t - 18\bigr).
  \label{eq:reward_pen}
\end{equation}
三项分别对应
（1）按图规模$N_t$归一化的强制重调度代价，
消除图越大绝对重调度次数天然越多的偏差；
（2）超出截止时限$T_D$的平滑线性惩罚，
取代硬截断以维持梯度连续性；
（3）拓扑膨胀惩罚，防止智能体通过无限分裂膨胀图规模
（阈值设为18，低于$N_{\max}=20$，为正常整形保留空间）。

% ----------------------------------------------------------------
\subsection{双图GNN架构}
\label{subsec:gnn}

状态$s_t$中的双图结构由三个神经网络子模块协同编码：
任务图节点编码器、资源图节点编码器，
以及将两图表示空间桥接起来的跨图注意力模块。
整体网络对应 \texttt{models/gnn\_policy.py}
中的 \texttt{NodeLevelActorCritic} 类。

\subsubsection*{节点级GNN编码器}

两个独立的 \texttt{\_GNNNodeEncoder}
（\texttt{gnn\_policy.py:276--333}）
分别对任务图与资源图进行编码。
编码器采用$L=2$层GraphSAGE~\cite{hamilton2017inductive}，
隐层维度$d_h = 64$，
最终经线性投影层映射至$d = 32$维嵌入空间。
与图级分类任务不同，编码器\emph{不执行}全局均值池化，
直接输出逐节点嵌入矩阵：
\begin{align}
  \mathbf{H}^T &= \textsc{GNN}_\tau\!\bigl(\mathbf{X}^T,\,\mathbf{E}^T\bigr)
    \in \mathbb{R}^{B \times N_{\max} \times d},
  \label{eq:task_enc}\\[3pt]
  \mathbf{H}^R &= \textsc{GNN}_r\!\bigl(\mathbf{X}^R,\,\mathbf{E}^R\bigr)
    \in \mathbb{R}^{B \times M \times d}.
  \label{eq:res_enc}
\end{align}
填充节点的零特征行在编码后被掩码清零，
防止其通过聚合操作污染有效节点的表示
（\texttt{gnn\_policy.py:463}）。

\subsubsection*{跨图注意力机制}

\texttt{CrossGraphAttention}（\texttt{gnn\_policy.py:340--380}）
实现从资源图到任务图的单头交叉注意力：
对每个任务节点$v_i$，
以其嵌入为查询，以全体无人机嵌入为键值对，
计算一个资源感知的上下文向量。
具体地，查询、键、值矩阵分别为
\begin{align}
  \mathbf{Q} &= \mathbf{H}^T\mathbf{W}_Q
    \in \mathbb{R}^{B \times N_{\max} \times d},
  \label{eq:cross_Q}\\
  \mathbf{K} &= \mathbf{H}^R\mathbf{W}_K
    \in \mathbb{R}^{B \times M \times d},
  \label{eq:cross_K}\\
  \mathbf{V} &= \mathbf{H}^R\mathbf{W}_V
    \in \mathbb{R}^{B \times M \times d},
  \label{eq:cross_V}
\end{align}
注意力权重矩阵为
\begin{equation}
  \boldsymbol{\Alpha} =
    \mathrm{softmax}\!\left(
      \frac{\mathbf{Q}\mathbf{K}^\top}{\sqrt{d}}
    \right) \in \mathbb{R}^{B \times N_{\max} \times M},
  \label{eq:attn_weights}
\end{equation}
其中$\Alpha_{b,i,k}$可被解释为
任务节点$v_i$对无人机$u_k$的软归属权重，
反映当前资源状态下$v_i$最可能被调度至哪台无人机。
资源增强的任务节点嵌入通过残差连接得到：
\begin{equation}
  \tilde{\mathbf{H}}^T =
    \mathbf{H}^T +
    \textsc{OutProj}\!\bigl(\boldsymbol{\Alpha}\mathbf{V}\bigr)
    \in \mathbb{R}^{B \times N_{\max} \times d}.
  \label{eq:enriched}
\end{equation}
注意力权重矩阵$\boldsymbol{\Alpha}$被完整保留
并传递至后续的合并评分头，
作为任务节点对之间链路质量的代理信号
（详见第~\ref{subsec:ppo}节）。

\subsubsection*{全局状态聚合}

全局状态向量$\mathbf{g}$汇聚双图信息与反馈信号，
用于驱动值函数头与空操作评分头
（\texttt{\_encode()} 函数，\texttt{gnn\_policy.py:458--494}）：
\begin{equation}
  \mathbf{g} = \Bigl[
    \bar{\mathbf{h}}^T_{\text{real}}\;\Big\|\;
    \bar{\mathbf{h}}^R\;\Big\|\;
    \mathbf{b}_{\text{res}}\;\Big\|\;
    \mathbf{f}(t)
  \Bigr] \in \mathbb{R}^{76},
  \label{eq:global_state}
\end{equation}
其中$\bar{\mathbf{h}}^T_{\text{real}}$为对真实（非填充）任务节点
求掩码均值所得的$d$维向量，
$\bar{\mathbf{h}}^R$为资源节点的均值嵌入，
$\mathbf{f}(t) \in \mathbb{R}^6$为调度反馈向量。
\emph{资源瓶颈向量}定义为
\begin{equation}
  \mathbf{b}_{\text{res}} =
  \Bigl[
    \min_k \mathbf{x}_k^R(t)\;\Big\|\;
    \frac{1}{M}\sum_k \mathbf{x}_k^R(t)
  \Bigr] \in \mathbb{R}^6,
  \label{eq:res_bottleneck}
\end{equation}
其最小值分量显式保留了机队中最薄弱无人机的状态，
防止均值聚合将单点故障信号淹没，
从而使值函数对个别无人机失效场景保持敏感
（\texttt{gnn\_policy.py:474--480}）。

% ----------------------------------------------------------------
\subsection{基于PPO的决策模块}
\label{subsec:ppo}

\subsubsection*{三头式动作Logit生成}

RTGS策略网络对式~\eqref{eq:action_dim}中的三类动作
分别使用独立评分头产生logit，
从而捕获不同粒度的动作语义
（\texttt{NodeLevelActorCritic.forward()}，
\texttt{gnn\_policy.py:499--536}）。

\smallskip
\noindent\textbf{空操作头}
以全局状态向量$\mathbf{g}$为输入，
输出单个标量logit：
\begin{equation}
  z_{\text{noop}} = \mathbf{w}_{\text{noop}}^\top \mathbf{g}
    \in \mathbb{R}.
  \label{eq:noop_logit}
\end{equation}

\smallskip
\noindent\textbf{分裂头}
以资源增强节点嵌入$\tilde{\mathbf{H}}^T$为输入，
对每个任务节点独立产生评分：
\begin{equation}
  \mathbf{z}_{\text{split}} =
    \tilde{\mathbf{H}}^T\,\mathbf{w}_{\text{split}}
    \in \mathbb{R}^{N_{\max}}.
  \label{eq:split_logit}
\end{equation}
由于$\tilde{\mathbf{H}}^T$已通过跨图注意力编码了任务节点
与各无人机的软归属关系，
智能体得以学习优先选择
由高CPU容量无人机承接的节点实施分裂。

\smallskip
\noindent\textbf{合并头}
为每个候选节点对$(v_i, v_j)$引入
\emph{链路质量信号}：
\begin{equation}
  q_{ij} =
    \boldsymbol{\Alpha}_i \cdot \boldsymbol{\Alpha}_j
  = \sum_{k=1}^{M} \Alpha_{ik}\,\Alpha_{jk} \in [0, 1].
  \label{eq:link_quality}
\end{equation}
$q_{ij}$高意味着$v_i$与$v_j$均主要关注同一台无人机，
两者可被同机执行而无需跨UAV传输，合并收益较低；
$q_{ij}$低则意味着两节点被分配至链路质量可能较差的不同无人机，
合并操作可消除该跨链路数据传输，从而降低端到端延迟。
合并logit由两层MLP计算：
\begin{equation}
  z_{ij}^{\text{merge}} =
    \textsc{MLP}_{\text{merge}}\!\bigl(
      [\,\tilde{\mathbf{h}}_i\;\|\;\tilde{\mathbf{h}}_j\;\|\;q_{ij}\,]
    \bigr) \in \mathbb{R},
  \label{eq:merge_logit}
\end{equation}
MLP输入维度为$2d + 1 = 65$
（\texttt{gnn\_policy.py:445--448}，
\texttt{merge\_head} 的结构为
$65 \to d \to 1$，中间采用ReLU激活）。

三组logit在动作维度上拼接，构成完整的动作评分向量：
\begin{equation}
  \mathbf{z} =
    \bigl[z_{\text{noop}}\;\|\;
          \mathbf{z}_{\text{split}}\;\|\;
          \mathbf{z}_{\text{merge}}\bigr]
    \in \mathbb{R}^{211}.
  \label{eq:all_logits}
\end{equation}

\subsubsection*{动作掩码}

为防止策略输出结构非法的图操作，
在Categorical采样前将无效动作的logit置为$-\infty$
（\texttt{compute\_action\_mask()}，
\texttt{gnn\_policy.py:539--584}）。
合法性规则如下：
\begin{itemize}[leftmargin=*, itemsep=2pt]
  \item \textbf{Noop}：恒为合法。
  \item \textbf{Split$(v_i)$}：当且仅当
        $v_i$为真实（非填充）节点、
        $N_t < N_{\max}$（图未满）
        且机队至少有一台无人机具有可用算力
        $\bigl(\max_k \tilde{\phi}_k(t) > 0\bigr)$时合法。
  \item \textbf{Merge$(v_i, v_j)$}：当且仅当
        $v_i$与$v_j$均为真实叶节点
        （在当前DAG中出度为零）时合法。
\end{itemize}
将布尔掩码记为$\mathbf{m} \in \{0,1\}^{211}$，
动作采样由掩码后的分布完成：
\begin{equation}
  a_t \sim \mathrm{Categorical}\!\bigl(
    \mathrm{softmax}(\mathbf{z} + \log \mathbf{m})
  \bigr),
  \label{eq:masked_sample}
\end{equation}
其中$\log 0 = -\infty$对应被屏蔽的无效动作。
掩码在策略回滚收集与PPO更新两个阶段均一致应用，
确保旧策略与新策略的对数概率计算在相同支撑集上进行
（\texttt{main.py:123--126}）。

\subsubsection*{PPO训练细节}

策略网络使用近端策略优化（PPO）算法~\cite{schulman2017proximal}
进行端到端训练
（\texttt{main.py:93--143}，\texttt{ppo\_update()} 函数）。
智能体首先收集$T_r = 2048$步滚动轨迹（约41个完整回合），
再对策略执行$K = 4$个更新轮次，
每轮次使用随机划分的小批量（大小$B_{\text{mb}} = 128$）
进行梯度更新。

复合损失函数为
\begin{equation}
  \mathcal{L} =
    \mathcal{L}_{\text{actor}}
    + c_v\,\mathcal{L}_{\text{critic}}
    - c_H\,\mathcal{H}(\pi),
  \label{eq:ppo_loss}
\end{equation}
其中actor损失采用PPO截断代理目标：
\begin{equation}
  \mathcal{L}_{\text{actor}} =
    -\,\mathbb{E}_t\!\left[
      \min\!\Bigl(
        \rho_t\,\hat{A}_t,\;
        \mathrm{clip}(\rho_t,\,1-\epsilon,\,1+\epsilon)\,\hat{A}_t
      \Bigr)
    \right],
  \label{eq:ppo_actor}
\end{equation}
$\rho_t = \pi_\theta(a_t|s_t)/\pi_{\theta_{\mathrm{old}}}(a_t|s_t)$
为新旧策略比率，
$\hat{A}_t$为归一化优势估计
$\bigl(\hat{A}_t \leftarrow (A_t - \mu_A)/(\sigma_A + \epsilon_0)\bigr)$，
截断系数$\epsilon = 0.2$；
critic损失为均方误差：
\begin{equation}
  \mathcal{L}_{\text{critic}} =
    \mathbb{E}_t\!\bigl[
      \bigl(V_\theta(s_t) - R_t\bigr)^2
    \bigr],
  \label{eq:ppo_critic}
\end{equation}
其中$R_t = \sum_{l \geq 0} \gamma^l r_{t+l}$为折扣累积回报
（折扣因子$\gamma = 0.99$）。
权重系数为$c_v = 0.5$、$c_H = 0.05$；
熵正则项$\mathcal{H}(\pi)$在211维动作空间中
显式鼓励探索，防止策略过早收敛至局部最优。
所有参数梯度以$\ell_2$范数上限$0.5$裁剪，
以保证大动作空间下的训练稳定性（\texttt{main.py:141}）。
总训练步数为$2\times10^6$步。

% ----------------------------------------------------------------
\subsection{RTGS完整决策流程}
\label{subsec:rtgs_flow}

每个决策时刻$t$的完整执行序列如下，
对应 \texttt{env/uav\_env.py:248--308}
中的 \texttt{step()} 方法：

\begin{enumerate}[label=\textbf{\arabic*.}, leftmargin=*, itemsep=4pt]

  \item \textbf{动作解码}：
        将平坦动作索引$a_t$通过
        \texttt{decode\_node\_action()} 解码为
        三元组$(op,\,\mathrm{pad}_i,\,\mathrm{pad}_j)$，
        再映射至图中实际节点ID $(v_i, v_j)$。

  \item \textbf{图整形}：
        调用 \texttt{AdaptiveTaskShaper.shape\_targeted()}，
        对$\mathcal{G}_T$执行
        \textsc{Split}$(v_i)$或\textsc{Merge}$(v_i, v_j)$算子，
        得到整形后的图$\mathcal{G}_T'$。

  \item \textbf{移动性感知贪心调度}：
        以当前带宽矩阵$\mathbf{B}(t)$
        （由式~\eqref{eq:bandwidth}计算）
        和有效CPU向量$\tilde{\boldsymbol{\phi}}(t)$
        （由式~\eqref{eq:cpu_derating}计算）为输入，
        执行 \texttt{greedy\_match()}，
        按拓扑序依次将每个任务分配至剩余CPU最多的无人机，
        返回makespan $T_{\mathrm{span}}(t)$、
        成功率$\sigma(t)$、
        重调度次数$\varrho(t)$等调度结果。

  \item \textbf{能量扣除}：
        按式~\eqref{eq:energy_depletion}
        根据本轮调度所产生的CPU消耗$\rho_k(t)$
        与传输量$\tau_k(t)$更新各无人机电量$\beta_k(t)$。

  \item \textbf{移动性推进}：
        按式~\eqref{eq:rwp}
        推进各活跃无人机的位置$\mathbf{p}_k(t+1)$，
        为下一时刻生成新的带宽矩阵$\mathbf{B}(t+1)$。

  \item \textbf{反馈更新与奖励计算}：
        构造新的调度反馈向量$\mathbf{f}(t)$，
        按式~\eqref{eq:reward}计算奖励$r_t$，
        并检查终止条件：
        若$T_{\mathrm{span}}(t) > 2\,T_D$则判定为灾难性失败并终止回合；
        若步数达到$T_{\max} = 50$则截断。

\end{enumerate}

上述流程构成RTGS的完整闭环：
每次图整形决策直接改变后续调度的任务粒度与并行度，
进而影响延迟与利用率，
并通过奖励信号反哺策略梯度更新，
使智能体逐步学会在动态异构UAV群环境下
做出最优的图拓扑自适应策略。
